\subsection{Existence for a positive exponent}

We prove part (1) of \cref{thm:als-prescribed-lp} by maximizing a weighted mass. Among the potentials $z\sim\varphi$ whose weighted measure $\mathrm{e}^{-\lambda z}\mathrm{MA}_\theta(z)$ lies below $\mu=f\,\mathrm{d}V$, we find one of maximal weighted mass; \cref{thm:relative-linfty} provides the compactness needed for this. If its weighted measure were strictly smaller than $\mu$, the addition theorem would produce a potential of larger weighted mass.

\Needspace{12\baselineskip}
\begin{proof}[of \cref{thm:als-prescribed-lp}, part (1)]
\emph{Step 1: a maximizing sequence.} Fix $\lambda>0$. By \cref{lem:relative-lp-supersolution}, there is $w\sim\varphi$ such that
\[
 \nu_0:=\mathrm{e}^{-\lambda w}\mathrm{MA}^{\mathrm{ext}}_\theta(w)
 \leq\mu=f\,\mathrm{d}V.
\]
Consider all potentials $z\sim\varphi$ for which
\[
 \nu_0\leq\nu_z:=\mathrm{e}^{-\lambda z}
                   \mathrm{MA}^{\mathrm{ext}}_\theta(z)\leq\mu,
\]
and let $B$ be the supremum of their weighted masses $\nu_z(X)$. This family is nonempty, since it contains $w$, and $B\leq\mu(X)<+\infty$. Choose $u_j$ in the family with $\nu_{u_j}(X)\to B$. Since $\nu_0\leq\nu_{u_j}$, \cref{cor:singular-exponential-order} gives $u_j\leq w$.

\emph{Step 2: uniform bounds.} We need a lower bound for the suprema of the $u_j$. Put $\varphi_\varepsilon=P_{\varepsilon\theta}(\varphi)\leq0$, and let $\psi$ be the bounded solution of
\[
 ((1-\varepsilon)\theta+\ddc\psi)^n=\mathrm{e}^{\lambda\psi}f\,\mathrm{d}V
\]
given by \cref{thm:big-bounded-exponential}. For $h=\psi+\varphi_\varepsilon$, positivity of the splitting $\theta_h=(\varepsilon\theta+\ddc\varphi_\varepsilon)+((1-\varepsilon)\theta+\ddc\psi)$ gives
\[
 \mathrm{e}^{-\lambda h}\mathrm{MA}^{\mathrm{ext}}_\theta(h)
 \geq \mathrm{e}^{-\lambda\varphi_\varepsilon}f\,\mathrm{d}V
 \geq\mu\geq\nu_{u_j}.
\]
Here $h\preceq u_j$, because $h\preceq\varphi$ and $u_j\sim\varphi$, and each $u_j$ admits a smaller-background envelope because $u_j\sim\varphi$. \Cref{cor:singular-exponential-order} therefore gives $u_j\geq h$. Consequently
\[
 \sup_Xh\leq M_j:=\sup_Xu_j\leq\sup_Xw,
\]
and the lower bound is finite.

We now apply \cref{thm:relative-linfty} to $v_j=u_j-M_j$. Their measures satisfy
\[
 \mathrm{MA}^{\mathrm{ext}}_\theta(v_j)
 =\mathrm{MA}^{\mathrm{ext}}_\theta(u_j)
 \leq \mathrm{e}^{\lambda\sup_Xw}f\,\mathrm{d}V,
\]
so $\varphi-C\leq v_j\leq\varphi$. Together with the bounds for $M_j$, this shows that all $u_j$ lie within a fixed bounded distance of $\varphi$.

\emph{Step 3: the limit.} Extract an $L^1$ limit $u$. By \cref{cor:common-barrier-continuity}, applied with the dominating measure $\mathrm{e}^{\lambda\sup_Xw}\mu$, the convergence holds in capacity, $u\sim\varphi$, and
\[
 \mathrm{MA}^{\mathrm{ext}}_\theta(u_j)
 \rightharpoonup\mathrm{MA}^{\mathrm{ext}}_\theta(u).
\]
After a further extraction we may also assume $\nu_{u_j}\rightharpoonup\nu$. Since $\nu_0\leq\nu_{u_j}\leq\mu$, the limit satisfies the same inequalities. We claim that
\[
 \mathrm{e}^{\lambda u_j}\nu_{u_j}\rightharpoonup \mathrm{e}^{\lambda u}\nu.
\]
Indeed, $\mathrm{e}^{\lambda u_j}\to \mathrm{e}^{\lambda u}$ in $L^1(\mu)$ by \cref{lem:fixed-npp-integrals}, and all the measures $\nu_{u_j}$ are at most $\mu$; this controls the error caused by changing the weight. The fixed bounded weight $\mathrm{e}^{\lambda u}$ can be approximated in $L^1(\mu)$ by continuous functions, and since both $\nu_{u_j}$ and $\nu$ are dominated by $\mu$, this approximation is uniform in the integrals. The weak convergence $\nu_{u_j}\rightharpoonup\nu$ then gives the claim, and the same argument applies after multiplication by any continuous test function.

\Needspace{7\baselineskip}
Since $\mathrm{e}^{\lambda u_j}\nu_{u_j}=\mathrm{MA}^{\mathrm{ext}}_\theta(u_j)$, it follows that
\[
 \mathrm{MA}^{\mathrm{ext}}_\theta(u)=\mathrm{e}^{\lambda u}\nu,
 \qquad \mathrm{e}^{-\lambda u}\mathrm{MA}^{\mathrm{ext}}_\theta(u)=\nu,
 \qquad\nu(X)=B.
\]
The second identity holds off the pole set of $u$, which has zero $\mu$-measure. Hence $u$ belongs to the family and attains the supremum $B$. Write $\nu=g\,\mathrm{d}V$, with $0\leq g\leq f$.

\emph{Step 4: no deficit.} Suppose that $g\neq f$ in $L^1(\mathrm{d}V)$. Applying \cref{lem:relative-lp-supersolution} to $f-g$, we find $v\sim\varphi$ with
\[
 \nu_v:=\mathrm{e}^{-\lambda v}\mathrm{MA}^{\mathrm{ext}}_\theta(v)
 \leq(f-g)\,\mathrm{d}V.
\]
This measure is nonzero. In fact, no potential $v$ admitting a smaller-background envelope has zero Monge--Amp\`ere measure, since \cref{thm:singular-strict-domination}, applied with coefficient zero to $v$ and $v+1$, would give $v\geq v+1$. Thus $\nu_v(X)>0$.

The sum $\nu+\nu_v$ is finite, so \cref{thm:weighted-measure-addition} gives $a\sim\varphi$ with
\[
 \mathrm{e}^{-\lambda a}\mathrm{MA}^{\mathrm{ext}}_\theta(a)=\nu+\nu_v.
\]
This measure still lies between $\nu_0$ and $\mu$, but its mass is $B+\nu_v(X)>B$, a contradiction. Therefore $g=f$, and $u$ solves the required equation. Uniqueness in the prescribed ordinary type follows from \cref{cor:singular-exponential-order}.
\end{proof}

\subsection{The normalized equation}

For the normalized equation we let the positive exponent tend to zero. Bounded comparison controls the normalizing constants from above, while the lower density estimate of \cref{lem:relative-density-lower-bound} keeps them away from zero. The common barriers then give convergence of the Monge--Amp\`ere measures and keep the limit in the prescribed type.

\Needspace{8\baselineskip}
\begin{proof}[of \cref{thm:als-prescribed-lp}, part (2)]
For each $j\geq1$, part (1) gives $u_j\sim\varphi$ with
\[
 \mathrm{MA}^{\mathrm{ext}}_\theta(u_j)=\mathrm{e}^{u_j/j}f\,\mathrm{d}V.
\]
Write $M_j=\sup_Xu_j$, $v_j=u_j-M_j$ and $c_j=\mathrm{e}^{M_j/j}$. Let $h$ be a bounded normalized solution of $\mathrm{MA}_\theta(h)=\alpha f\,\mathrm{d}V$, with $\alpha>0$, given by \cref{thm:big-bounded-normalized}.

\Needspace{5\baselineskip}
For $H_j=h+\|h\|_\infty+j\log\alpha$ we have
\[
 \mathrm{e}^{-H_j/j}\mathrm{MA}_\theta(H_j)
 =\mathrm{e}^{-(h+\|h\|_\infty)/j}f\,\mathrm{d}V
 \leq f\,\mathrm{d}V
 =\mathrm{e}^{-u_j/j}\mathrm{MA}^{\mathrm{ext}}_\theta(u_j).
\]
The bounded function $H_j$ admits a smaller-background envelope, and $u_j\preceq H_j$, so \cref{cor:singular-exponential-order} gives $u_j\leq H_j$. In particular,
\[
 0<c_j\leq\alpha\exp(2\|h\|_\infty/j).
\]

The lower density bound of \cref{lem:relative-density-lower-bound}, applied to $v_j$ with $q=p$, gives
\[
 b_p\leq\|c_j\mathrm{e}^{v_j/j}f\|_p\leq c_j\|f\|_p,
\]
so the $c_j$ are also bounded away from zero.

Since $\sup_Xv_j=0$ and $\mathrm{MA}^{\mathrm{ext}}_\theta(v_j)=c_j\mathrm{e}^{v_j/j}f\,\mathrm{d}V\leq C f\,\mathrm{d}V$, \cref{thm:relative-linfty} gives $\varphi-C'\leq v_j\leq\varphi$. Extract a subsequence with $v_j\to v$ in $L^1(X)$ and $c_j\to c>0$. By \cref{cor:common-barrier-continuity}, the convergence holds in capacity, $v\sim\varphi$, and the Monge--Amp\`ere measures converge weakly. Meanwhile, \cref{lem:normalized-exponential-lp} gives
\[
 \int_X|\mathrm{e}^{v_j/j}-1|f\,\mathrm{d}V\longrightarrow0.
\]
Thus the right-hand sides converge in total variation to $cf\,\mathrm{d}V$, and $\mathrm{MA}^{\mathrm{ext}}_\theta(v)=cf\,\mathrm{d}V$.

For the uniqueness of the constant, suppose that $a,b\sim\varphi$ solve the normalized equation with constants $c_a<c_b$. For every real $t$,
\[
 \mathrm{MA}^{\mathrm{ext}}_\theta(a)
 =\frac{c_a}{c_b}\mathrm{MA}^{\mathrm{ext}}_\theta(b+t),
\]
so \cref{thm:singular-strict-domination} gives $a\geq b+t$. This is impossible as $t\to+\infty$, at any point where both potentials are finite. Interchanging the roles of the two constants proves that $c$ is unique.
\end{proof}
